\documentclass[12pt]{article}
\usepackage[T1]{fontenc}
\usepackage[utf8]{inputenc}
\usepackage{lmodern}
\usepackage{textcomp}
\usepackage{enumitem}
\usepackage{geometry}
\usepackage{graphicx}
\usepackage{amsmath, amssymb}
\geometry{margin=1in}
\begin{document}
\begin{center}
Philosophy - Graduate Level Assessment 22 \
Total Marks: 40 \
Answer all questions.
\end{center}

\section*{Multiple Choice (10 marks)}
\begin{enumerate}[label=\arabic*.]
\item Hobbes describes felicity as:
\begin{enumerate}[label=(\alph*)]
    \item the absence of desire.
    \item a state of constant dissatisfaction.
    \item a state of constant change.
    \item the absence of progress.
    \item a state of satisfaction with one's current state.
\end{enumerate}
\item Ashford's article is meant to address a particular paralysis in the face of
\begin{enumerate}[label=(\alph*)]
    \item the daunting task of solving worldwide economic imbalance.
    \item the impossibility of meeting everyone's basic needs.
    \item having to give up so much of our own wealth in light of Singer's arguments.
    \item having to choose between addressing immediate harm and addressing structural injustice.
    \item reconciling conflict moral theories.
\end{enumerate}
\item Craig says an actually infinite number of things \_\_\_\_\_.
\begin{enumerate}[label=(\alph*)]
    \item cannot exist
    \item can be physically observed
    \item can be counted one by one
    \item is a concept beyond human comprehension
    \item only exists in mathematics
\end{enumerate}
\item Naturalists who concentrated on natural elements and processes are associated with which of the following?
\begin{enumerate}[label=(\alph*)]
    \item Mengzi
    \item Zoroastrianism
    \item Zen Buddhism
    \item Shintoism
    \item Vedic Philosophy
\end{enumerate}
\item Lukianoff and Haidt argue that American colleges and universities now encourage
\begin{enumerate}[label=(\alph*)]
    \item emotional reasoning.
    \item critical reasoning.
    \item reflective reasoning.
    \item all of the above.
\end{enumerate}
\end{enumerate}

\section*{True / False (10 marks)}
\textit{Answer True or False}
\begin{enumerate}[label=\arabic*.]
\setcounter{enumi}{5}
\item True or False: Sinnott-Armstrong mainly investigates the moral responsibilities of technological companies.
\item True or False: Baxter would likely argue that the policy is morally unproblematic since it does not adversely affect human beings.
\item True or False: According to Singer, the conclusions in 'All Animals Are Equal' stem primarily from the principal principle of moral consideration.
\item True or False: Gauthier argues that moral agreements that are equally favorable to all parties are desirable because they satisfy our desire for fairness.
\item True or False: The Ideal Moral Code theory posits that moral obligations are contingent upon individual interpretations of ethical rules.
\end{enumerate}

\section*{Long Form Response (20 marks)}
\textit{Answer each question with a clear and complete explanation.}
\begin{enumerate}[label=\arabic*.]
\setcounter{enumi}{10}
\item Analyze the logical structure of the following statement: "Either England's importing beef is a sufficient condition for France's subsidizing agriculture, or China doesn't promote human rights when and only when South Africa supplies diamonds." How can this be symbolized in propositional logic, and what implications does this have for understanding conditional relationships in ethical discourse?
\item Discuss the concept of moral cosmopolitanism, analyzing the obligations it places on individuals to assist those in need and promote basic human rights.
\end{enumerate}

\vfill
\noindent\textit{End of Paper}
\end{document}